\subsection{対称性と漸化式}

$u=1-x$と置くと二つの指数が入れ替わるため
\begin{equation}\label{式:ベータ対称性}
 B(p,q)=B(q,p)
\end{equation}
である。広義積分の場合も，端点を避けた区間で置換し，各端点の極限を取れば同じ結論になる。

また$x+(1-x)=1$を掛ければ
\begin{equation}\label{式:ベータ和の関係}
 B(p+1,q)+B(p,q+1)=B(p,q).
\end{equation}
もう一つの関係を得るために$x^p(1-x)^q$を微分する。
\[
 \frac{\dd}{\dd x}\{x^p(1-x)^q\}
 =p x^{p-1}(1-x)^q-qx^p(1-x)^{q-1}.
\]
これを$\varepsilon$から$1-\delta$まで積分し，両端をそれぞれ極限へ近づける。$p,q>0$なので$x^p(1-x)^q$は両端で$0$へ近づき，右辺の二つの積分はすでに収束すると分かっている。よって
\[
 pB(p,q+1)=qB(p+1,q).
\]
式\eqref{式:ベータ和の関係}と連立すると
\begin{equation}\label{式:ベータ漸化式}
 B(p+1,q)=\frac{p}{p+q}B(p,q),\qquad
 B(p,q+1)=\frac{q}{p+q}B(p,q)
\end{equation}
を得る。この漸化式は整数以外にも成り立つ。整数のときには，$B(1,1)=1$から出発して階乗の公式を再び導くこともできる。

